%==============================================================================
%  DRAFT -- ICRA / RA-L submission.
%  Occlusion-aware marker-free human following for MAVs in forests.
%
%  THESIS (from the annotated bibliography, 2026-09-10):
%    Every prior aerial-tracking system treats target loss as ONE undifferentiated
%    event and answers it with ONE policy: extrapolate and keep flying.  We
%    classify the CAUSE of loss against the onboard-built 3-D map -- occluded by
%    known geometry vs. genuinely lost -- and condition planner behaviour on that
%    classification.  The geometry needed is already computed for free as a
%    by-product of obstacle avoidance.
%
%  STATUS 2026-09-10: skeleton + method + system characterization.  NO FLIGHT
%  RESULTS.  Conventions:
%      \tbd{...}   red    -- not written or not measured yet
%      \bench{...} blue   -- a REAL measurement, but bench/handheld, NEVER in flight
%      \note{...}  gray   -- drafting instruction to ourselves; delete before submit
%
%  BUILD: Overleaf / pdfLaTeX.  Never compiled -- there is no TeX toolchain on the
%  aircraft's Jetson.  Expect first-build fixes.
%
%  CLASS: IEEEtran conference.  For ICRA's own template swap in
%      \documentclass[letterpaper, 10pt, conference]{ieeeconf}
%  and drop \IEEEoverridecommandlockouts.  ICRA = 6 pages + up to 2 paid.
%  RA-L = 8 pages.  This draft is written to the 8-page shape; cut Sec. VIII
%  first if targeting ICRA at 6.
%==============================================================================
\documentclass[conference]{IEEEtran}
\IEEEoverridecommandlockouts

\usepackage{cite}
\usepackage{siunitx}
\usepackage{amsmath,amssymb,amsfonts}
\usepackage{algorithm}
\usepackage{algorithmic}
\usepackage{graphicx}
\usepackage{booktabs}
\usepackage{xcolor}
\usepackage{multirow}
\usepackage[hidelinks]{hyperref}
\usepackage{url}

% --- drafting aids.  Strip the color (or the macros) before submitting. -------
\newcommand{\tbd}[1]{\textcolor{red}{\textbf{[TBD: #1]}}}
\newcommand{\bench}[1]{\textcolor{blue}{#1}}
\newcommand{\note}[1]{\textcolor{gray}{\textit{(#1)}}}
% Shorthands for the two loss classes, so the naming stays consistent.
\newcommand{\occluded}{\textsc{Occluded}}
\newcommand{\lost}{\textsc{Lost}}
% -----------------------------------------------------------------------------

\begin{document}

% ALTERNATE TITLES, pick one:
%  - "Knowing Why You Lost Them: Map-Conditioned Occlusion Reasoning for
%     Marker-Free Human Following in Forests"
%  - "Occlusion-Aware Human Following for MAVs in Cluttered Natural Environments"
%  - "Classifying Target Loss for Robust Aerial Human Following"
\title{Knowing \emph{Why} You Lost Them: Map-Conditioned Occlusion\\
Reasoning for Marker-Free Human Following in Forests}

\author{%
\IEEEauthorblockN{Kevin Biju Mathew}
\IEEEauthorblockA{Robotics / Interdisciplinary Engineering\\
Purdue University\\
West Lafayette, IN, USA\\
\tbd{email}}
\and
\IEEEauthorblockN{Shrikrishna Bhagirath Rajule}
\IEEEauthorblockA{Robotics / Interdisciplinary Engineering\\
Purdue University\\
West Lafayette, IN, USA\\
\tbd{email}}
\and
\IEEEauthorblockN{Aniket Bera}
\IEEEauthorblockA{Department of Computer Science\\
Purdue University\\
West Lafayette, IN, USA\\
\tbd{email}}
}

\maketitle

%==============================================================================
\begin{abstract}
\note{Write LAST, 150--200 words. Beats: (1) aerial following in forests, (2)
everyone treats loss as one event, (3) we classify the cause against the map,
(4) result.}
Autonomous aerial human following fails in forests, and it fails in a specific
way: the target disappears behind a trunk. Every existing occlusion-aware aerial
tracker responds to target loss with a single policy --- extrapolate the last
motion estimate and keep flying --- because loss is treated as one
undifferentiated event. We argue that it is two events with opposite correct
responses. A subject who has stepped behind a \SI{}{}\tbd{width}\,m trunk at
\SI{1.4}{\meter\per\second} will re-emerge on a predictable side in under a
second, and the right behaviour is to hold standoff and pre-position for that
edge; a subject who has genuinely left the sensor's reach requires a search. We
show that the geometry needed to tell these apart is already present in the
occupancy map an obstacle-avoiding MAV builds for free, and we present a
following system that queries it: a line-of-sight test against known geometry
classifies each loss as \occluded{} or \lost{}, a reappearance belief is
propagated over the occluder's silhouette, and the planner is conditioned on the
resulting class. The system runs entirely onboard \tbd{platform one-liner},
tracks an \emph{unmarked} human --- unlike prior LiDAR aerial tracking, which
instruments its targets with retroreflective tape --- and we evaluate it in
\tbd{N} real forest trials, reducing re-acquisition time by \tbd{X}\,\% against
\tbd{baseline}.
\end{abstract}

\begin{IEEEkeywords}
Aerial systems: perception and autonomy; visual tracking; motion and path
planning; field robotics; human detection and tracking.
\end{IEEEkeywords}

%==============================================================================
\section{Introduction}
\label{sec:intro}

\note{~1 column + teaser. Beats in order:}

\begin{enumerate}
  \item \textbf{The task.} An MAV that keeps a designated person in view through
        natural clutter --- search-and-rescue escort, forestry and field work
        support, cinematography under canopy.
  \item \textbf{The forest makes it hard in one particular way.} Trunks are
        narrow, dense, and everywhere; occlusion is not an edge case but the
        steady state. A commercial visual-tracking drone fails at the woodland
        boundary~\cite{ren2025super}.
  \item \textbf{What everyone does about it.} The occlusion-aware aerial
        tracking literature is mature and
        capable~\cite{ji2022elastic,wang2021visibility,svpto2023,gao2024adaptive},
        and the most recent work adds explicit
        re-acquisition~\cite{lin2026evatracker}. But target loss is handled as a
        single event class: predict forward and continue.
  \item \textbf{Why that is the wrong abstraction.} State the trunk example
        numerically. Two losses, same sensor signature, opposite correct
        responses. An unconditional extrapolation policy is right for one and
        wasteful-or-dangerous for the other.
  \item \textbf{The observation that makes it cheap.} The information required
        to separate them --- is the predicted target position behind
        \emph{known} occupied geometry, as seen from the current viewpoint? ---
        is a ray query into the occupancy map the planner already
        maintains~\cite{ren2023rogmap}. No new sensor, no new map, no learned
        occlusion model.
  \item \textbf{Contributions.}
\end{enumerate}

\noindent\textbf{Contributions.}
\begin{itemize}
  \item \textbf{Loss-cause classification.} A line-of-sight test against the
        onboard-built occupancy map that labels every target-loss event
        \occluded{} or \lost{}, using geometry already computed for obstacle
        avoidance (Sec.~\ref{sec:classify}).
  \item \textbf{A reappearance belief, not a point.} Under \occluded{}, we
        propagate the predicted target state over the occluder's silhouette to
        produce a distribution over \emph{where and when} the subject re-emerges,
        and condition the viewpoint planner on it (Sec.~\ref{sec:reappear}).
  \item \textbf{Marker-free LiDAR-and-RGBD human following.} Prior LiDAR aerial
        tracking, including the system we build
        on~\cite{ren2025super} and its swarm
        extension~\cite{yin2023swarm}, instruments the target with
        retroreflective tape. We track an unmarked human.
  \item \textbf{Forest evaluation with an uninstrumented human subject}, against
        three baselines including a re-acquisition-capable prior
        system~\cite{lin2026evatracker}, plus an offline prediction study under
        induced missing observations~\cite{chib2024trajimpute} that isolates the
        predictor from the flight system (Sec.~\ref{sec:results}).
  \item \textbf{A layered safety architecture} in which the follow behaviour
        cannot outlive its own evidence, which is what makes an autonomous
        follower safe to fly near the person it is following
        (Sec.~\ref{sec:safety}).
\end{itemize}

\begin{figure}[t]
  \centering
  \framebox[0.95\columnwidth]{\rule{0pt}{4.5cm}\ \tbd{TEASER. Best version:
  three panels of the same instant --- (a) onboard camera, subject vanishing
  behind a trunk; (b) the occupancy map with the LoS ray blocked and the
  reappearance region shaded on the trunk's left edge; (c) the flown trajectory
  pre-positioning for that edge. The figure IS the paper's argument.}}
  \caption{\tbd{caption}}
  \label{fig:teaser}
\end{figure}

%==============================================================================
\section{Related Work}
\label{sec:related}

\note{~1 column, 3 subsections, then the table. Structure fixed --- each
subsection ends by naming the shared limitation.}

\subsection{Perception- and Visibility-Aware Aerial Tracking}
Aerial target tracking in clutter matured through target-informed kinodynamic
search~\cite{han2021fasttracker}, differentiable visibility
metrics folded into trajectory
optimization~\cite{wang2021visibility,jeon2019online,bonatti2018cinematographer},
and joint spatio-temporal optimization with analytic occlusion
costs~\cite{ji2022elastic}, consolidated in~\cite{gao2024adaptive}. Related
formulations adjust standoff distance and relative angle with environment
complexity~\cite{svpto2023}, optimize position and yaw
jointly~\cite{tordesillas2022panther}, or learn the occlusion
cost~\cite{masnavi2022visibility}. The most recent work adds explicit
re-acquisition: prediction-driven path generation continues from predicted target
positions after loss~\cite{lin2026evatracker}.
\emph{Shared limitation:} occlusion is treated either as a cost to \emph{avoid}
or as an event to \emph{extrapolate through}. No system asks why the target
disappeared. \note{This sentence is the paper's hinge --- make sure it is
defensible against all six citations above.}

\subsection{Human Trajectory Prediction Under Partial Observation}
Pedestrian forecasting is a mature field~\cite{rudenko2020survey}, from social
pooling~\cite{alahi2016sociallstm,gupta2018socialgan} through dynamics-constrained
graph models~\cite{salzmann2020trajectron} to low-rank descriptor
spaces~\cite{bae2023eigentrajectory} and diffusion
predictors~\cite{bae2024singulartrajectory,mao2023led}. Nearly all of it assumes
a \emph{complete} observation history; \cite{chib2024trajimpute} states the
assumption explicitly and provides the first benchmark for missing observations,
alongside joint imputation-and-prediction
approaches~\cite{xu2023uncovering,becker2021missformer}. Closest to our reasoning
is occlusion inference for partially observable driving
scenes~\cite{shi2024sceneinformer}, which jointly infers occluded agents and
predicts observed ones.
\emph{Shared limitation:} these methods are not coupled to a planner, and the
map-conditioned variants assume a prior HD map rather than an occupancy grid the
robot builds online.

\subsection{LiDAR MAV Navigation in Unstructured Environments}
Corridor-based high-speed planning~\cite{ren2022bubble,tordesillas2021faster},
ESDF-free local planning~\cite{zhou2020egoplanner}, perception-aware
replanning~\cite{zhou2021raptor}, and the MINCO trajectory
class~\cite{wang2022minco} underpin modern LiDAR MAV autonomy, with
tightly-coupled LiDAR-inertial odometry~\cite{xu2022fastlio2} and
planning-oriented occupancy maps~\cite{ren2023rogmap} as the substrate. Field
deployments now reach real forests~\cite{schedl2021forest,boreal2025}, and
learning-based flight has been demonstrated in the
wild~\cite{loquercio2021highspeed}.
\emph{Shared limitation:} the tracking modes of these systems use
\emph{instrumented} targets --- retroreflective markers --- and constant-velocity
motion models~\cite{ren2025super,yin2023swarm}.

\begin{table*}[t]
\centering
\caption{Positioning against the closest prior systems. \note{Reviewers in this
subfield expect this table. Fill every cell --- an empty cell reads as an
unexamined claim.}}
\label{tab:related}
\footnotesize
\begin{tabular}{@{}lllllll@{}}
\toprule
System & Occlusion handling & Target prediction & Loss-cause & Reappearance
& Target & Environment \\
       &                    &                   & classified? & belief? & marked? & \\
\midrule
Elastic Tracker~\cite{ji2022elastic}     & preventive cost      & intention-free & no & no & \tbd{} & obstacles \\
Wang et al.~\cite{wang2021visibility}    & differentiable cost  & \tbd{}         & no & no & \tbd{} & obstacles \\
SVPTO~\cite{svpto2023}                   & standoff modulation  & \tbd{}         & no & no & \tbd{} & obstacles \\
BPOD~\cite{bpod2023}                     & belief-space $P$(detect) & stochastic & no & \emph{implicit} & \tbd{} & clutter \\
Eva-Tracker~\cite{lin2026evatracker}     & FoV-ESDF             & analytic/geometric & \textbf{no} & no & \tbd{} & clutter \\
SUPER~\cite{ren2025super}                & none (safety only)   & constant velocity & no & no & \textbf{yes} & forest \\
\textbf{Ours}                            & map raycast + planner conditioning & CTRV\,$\oplus$\,learned & \textbf{yes} & \textbf{yes} & \textbf{no} & forest \\
\bottomrule
\end{tabular}
\end{table*}

\subsection{Explicit Differentiation}
\note{One paragraph, one sentence per Tier-0 system. A reviewer will check
exactly this. Drafted below --- sharpen once you have read all four in full.}

Against~\cite{lin2026evatracker}, which shares our hardware class and our
re-acquisition claim: their loss handling is unconditional prediction
carry-forward, so a target that walked out of range and a target standing behind
a trunk receive identical treatment, and their FoV-ESDF is a planner-side
visibility construct rather than a belief over re-emergence.
Against~\cite{ji2022elastic}: their occlusion cost is \emph{preventive}, with no
model of behaviour \emph{during} occlusion, and its stated intention-free design
is precisely the prior we argue is wrong for a human target in structured
clutter. Against~\cite{bpod2023}, the closest formalization of reasoning about
future visibility: their probability of detection is propagated over continuous
stochastic states, where ours is a discrete-event, map-conditioned classification
that changes the planner's \emph{mode} rather than weighting a cost.
Against~\cite{ren2025super}, our own base planner: its tracking mode uses a
retroreflective marker and a constant-velocity model, which makes its occlusion
story untested by construction --- the target is trivially detectable whenever it
is visible at all. \note{Also address \cite{bera2016glmp} or your own prior work
if relevant --- self-citation placement is your call.}

%==============================================================================
\section{Problem Formulation}
\label{sec:problem}

\note{New section --- the old draft had none, and this thesis needs one. Keep it
short and make the classification precise enough to implement from the paper.}

Let $\mathcal{M}_t$ be the onboard occupancy map at time $t$, assigning each
voxel one of $\{\textsc{Free}, \textsc{Occupied}, \textsc{Unknown}\}$, and let
$\mathbf{p}^{\text{mav}}_t$ be the vehicle position. A target track carries state
$\mathbf{x}_t = [\,\mathbf{p}_t, \mathbf{v}_t\,]$ and a predictor
$f(\cdot)$ producing $\hat{\mathbf{p}}_{t+\tau}$ for a horizon $\tau$.

At the first frame $t^\ast$ in which the tracker reports no detection for the
locked target, we evaluate the line segment
$\overline{\mathbf{p}^{\text{mav}}_{t^\ast}\,\hat{\mathbf{p}}_{t^\ast}}$ against
$\mathcal{M}_{t^\ast}$ and assign

\begin{equation}
c_{t^\ast} =
\begin{cases}
\occluded{} & \text{if the segment intersects \textsc{Occupied} voxels} \\
             & \text{before reaching } \hat{\mathbf{p}}_{t^\ast},\\[2pt]
\lost{}      & \text{otherwise.}
\end{cases}
\label{eq:classify}
\end{equation}

\tbd{This is the honest first version and it is too simple. Decide and write up:
(i) the \textsc{Unknown} case --- a ray through unmapped space is neither
occlusion nor loss, and in a forest it is common; propose a third class
\textsc{Unmapped} or fold it into \lost{} and say why. (ii) A single ray to a
point target is brittle; use the target's bounding volume and report the
\emph{fraction} of rays blocked, thresholded. (iii) Detector false negatives
produce loss events with a clear line of sight --- these correctly classify as
\lost{}, which is a feature, but say so. (iv) State the time budget: this must
run inside one \SI{8}{\hertz} (\SI{125}{\milli\second}) perception frame.}

Under $c=\occluded{}$ we additionally require the \emph{reappearance set}
$\mathcal{R} \subset \mathbb{R}^3$: positions on the boundary of the occluding
geometry at which the predicted target trajectory re-enters free, visible space,
with an associated time distribution (Sec.~\ref{sec:reappear}).

%==============================================================================
\section{System Overview}
\label{sec:system}

\subsection{Platform}
\begin{table}[t]
\centering
\caption{Onboard hardware. All compute figures in this paper
(Table~\ref{tab:compute}) were measured on this module.}
\label{tab:hardware}
\footnotesize
\begin{tabular}{@{}ll@{}}
\toprule
Component & Part / configuration \\
\midrule
Airframe        & Gypsy MK5, \tbd{mass, prop size, endurance} \\
LiDAR           & Livox MID-360, \SI{10}{\hertz} frames, \SI{200}{\hertz} IMU \\
Camera          & Orbbec Gemini 336L, active stereo, $1280\times720$ @ \SI{30}{\hertz} \\
Compute         & NVIDIA Jetson Orin NX (16\,GB), ROS~2 Humble \\
Autopilot       & ARK FPV, PX4 v1.17, STM32H7 \\
Companion link  & uXRCE-DDS over UART @ 3\,Mbaud \\
\bottomrule
\end{tabular}
\end{table}

\subsection{Architecture}
\begin{figure}[t]
  \centering
  \framebox[0.95\columnwidth]{\rule{0pt}{5.5cm}\ \tbd{ARCHITECTURE FIGURE. Must
  show: the map serving BOTH the planner and the loss classifier (that shared
  edge is the contribution); the two independent paths into the safety veto; and
  the goal-status return edge from the planner.}}
  \caption{\tbd{caption}}
  \label{fig:arch}
\end{figure}

Six stages, each publishing only while it holds a belief it can defend:
LiDAR-inertial state estimation; human detection and identity-persistent
tracking; trajectory prediction; \textbf{loss classification and reappearance
belief}; standoff viewpoint planning; and goal arbitration with an independent
human-proximity veto. Stage four is the contribution; the rest is engineering we
describe only to the extent that the contribution depends on it.

%==============================================================================
\section{Method}
\label{sec:method}

\subsection{State Estimation and Autopilot Integration}
\label{sec:state}
We use tightly-coupled iterated-EKF LiDAR-inertial odometry~\cite{xu2022fastlio2}
as the sole absolute position source, fused by the autopilot's estimator as an
external-vision measurement alongside optical-flow velocity and a downward
rangefinder; satellite navigation is disabled so reported behaviour is
attributable to the onboard estimate. Three interface details are load-bearing
and are, in our experience, where this integration is usually silently wrong:
the odometry is stamped at the instant the filter propagates state to, so the
estimator's measurement-delay parameter must be \emph{zero} or the latency is
double-counted; the odometry frame is the LiDAR's internal IMU, not the vehicle
body, so the declared lever arm must be verified against the constant offset
between estimator output and raw odometry \bench{(measured $(-0.014, +0.024,
+0.064)$\,m against a declared $(+0.011,-0.023,-0.070)$\,m)}; and mounting
\emph{tilt} is recoverable by fusing vision yaw only and differencing
gravity-referenced attitude from two independent IMUs \bench{(roll $-1.55^\circ$,
pitch $+0.09^\circ$)}, while mounting \emph{yaw} is unobservable from the
estimator that consumes it and must be established mechanically.
\note{Compress to 4--5 sentences if page-limited; the detail belongs in
Sec.~\ref{sec:lessons} anyway.}

\subsection{Detection, Tracking, and 3-D Localization}
\label{sec:perception}
Per colour frame: a single-stage detector restricted to the person class, with
identity association by a motion-and-appearance
tracker~\cite{aharon2022botsort,zhang2022bytetrack}. Each confirmed track is
localized in 3-D by deprojecting depth into an organized point array and reducing
the samples inside the track's box. World$\leftarrow$camera is taken at the
frame's \emph{own} timestamp, never the latest transform: with a
\SI{10}{\hertz} pose source and a \SI{30}{\hertz} camera the difference is up to
a full pose interval of position error on a moving vehicle, and it is invisible
in static testing.

\paragraph{Depth source} Two candidates are carried and selected at run time.
\bench{Indoors, the camera's active stereo yields \SI{88.4}{\percent} valid
pixels at \SI{30}{\hertz}; LiDAR points projected through the solved extrinsic
cover \SI{13.5}{\percent} of the camera field of view, with no returns in the
lower third owing to the sensor's ${\approx}-7^\circ$ vertical floor.} The camera
is the default. The caveat is structural rather than incidental --- it is an
\emph{active} stereo sensor whose projector competes with solar infrared, so the
indoor advantage is exactly what outdoor operation attacks, and canopy
illumination is neither indoor nor open-sky.
\tbd{Outdoor + under-canopy depth characterization. This experiment is designed,
not run, and a reviewer WILL ask. See also~\cite{canopy2026} on LiDAR return
behaviour through vegetation.}

\paragraph{Extrinsics} Targetless single-shot camera--LiDAR
calibration~\cite{koide2023calibration} with learned
correspondences~\cite{sarlin2020superglue}. \bench{8 capture pairs, validated at
$\Delta u = 0$\,px, $\Delta v = +1$\,px.} The calibration is bound to specific
sensor serial numbers and refuses to load against a mismatched pair.
\tbd{Motor-on re-validation --- solved with motors off; vibration shifts it.}

\subsection{Trajectory Prediction}
\label{sec:prediction}
We fuse a kinematic prior with a learned pedestrian model. The kinematic branch
is a constant-turn-rate-and-velocity
model~\cite{schubert2008comparison}; the learned branch is a low-rank trajectory
descriptor predictor~\cite{bae2023eigentrajectory} operating on an
$8$-step observation window with a $12$-step horizon, blended with weight
$\alpha$.

Two properties of this pairing matter for the contribution:

\begin{itemize}
  \item The learned predictor \textbf{assumes a complete observation
        history}~\cite{chib2024trajimpute}. Occlusion is precisely the condition
        that violates it, so the fusion weight must be a function of observation
        completeness rather than a constant. \tbd{state the actual schedule ---
        currently a fixed $\alpha$; make it depend on missed-frame count and say
        so, or justify the constant.}
  \item Under \occluded{}, the kinematic branch is the one that carries a
        physically meaningful prior through the gap, while the learned branch
        contributes multimodality about \emph{which way} the subject went around
        the occluder. \tbd{verify this is what actually happens; it is the
        claim the ablation in Sec.~\ref{sec:results} tests.}
\end{itemize}

\note{Honest status for the writing team: the learned branch ships DISABLED by
default and its checkpoint/config live outside version control. Either enable and
characterize it, or describe the system as CTRV with the learned branch as an
evaluated option --- but the paper must not imply it flew if it did not.}

\subsection{Loss-Cause Classification}
\label{sec:classify}
\note{THE CONTRIBUTION. This subsection deserves the most space, an algorithm
block, and a figure. Everything above is setup.}

On the first frame with no detection for the locked target, we evaluate
Eq.~\eqref{eq:classify} by casting into the planner's occupancy
map~\cite{ren2023rogmap}.

\begin{algorithm}[t]
\caption{\tbd{Loss classification and reappearance belief}}
\label{alg:classify}
\begin{algorithmic}[1]
\STATE \tbd{Write this properly once the implementation settles. Shape:}
\STATE \textbf{input:} map $\mathcal{M}$, MAV pose, last confirmed track, predictor $f$
\STATE $\hat{\mathbf{p}} \leftarrow f(\text{history}, \tau)$
\STATE cast $N$ rays from MAV to the predicted target volume
\STATE $\rho \leftarrow$ fraction blocked by \textsc{Occupied} before $\hat{\mathbf{p}}$
\STATE \textbf{if} $\rho > \rho_{\text{thresh}}$ \textbf{then} $c \leftarrow \occluded{}$
       \textbf{else} $c \leftarrow \lost{}$
\STATE \textbf{if} $c = \occluded{}$ \textbf{then} compute $\mathcal{R}$, the
       silhouette-boundary crossing set, with arrival-time distribution
\STATE \textbf{return} $c$, $\mathcal{R}$
\end{algorithmic}
\end{algorithm}

\paragraph{Why this is nearly free} The map is maintained for obstacle avoidance
regardless; the classification adds $N$ ray queries per lost frame.
\tbd{measure and report: $\mu$s per classification, and the worst case within a
\SI{8}{\hertz} budget. A cheap contribution is a feature --- quantify it.}

\paragraph{The dependency, and its one-line configuration trap}
The classifier inherits whatever the map can represent. If the map frees cells
only by raycasting and raycasting is off, the map never forgets: a walking person
is inserted as permanent occupancy, the subject stamps a wall behind themselves
along the corridor the vehicle must fly, and Eq.~\eqref{eq:classify} then queries
that stale occupancy and reports \occluded{} against the target's own history.
\tbd{Raycasting is now ENABLED in the flown follow configuration and its patch
(2026-09-12); it was previously disabled, inherited by copy from the bench
configuration. Still to measure before flying: its CPU cost, reported alongside
the classifier's own; clearing time on hardware; and trail false-\occluded{}. A
saturated cell needs 31 clean rays to leave \textsc{Occupied} and 74 to become
\textsc{Free} at $p_{\max}=0.95$, reading \textsc{Unknown} in between, and a subject
walking directly away leaves that trail on the line of sight. If the locked
target's points are filtered from the map input (plan M5.3), state it here.}
\note{Worth a sentence in the paper regardless of our own bug: the dependency is
real and general --- any map that cannot clear a moving occupant will corrupt this
classification, so it is a stated precondition of the method, not an
implementation detail.}
\note{Second precondition, found 2026-09-14 on a bench replay (docs/m6-loss-oracle.md): the
LiDAR's vertical field of view bounds what the map can know. The MID360 looks no lower than
$7^\circ$ below its horizon, so at range $d$ nothing below $z_{\text{sensor}} - d\tan 7^\circ$
is ever observed --- \SI{0.5}{\meter} under the sensor at \SI{4}{\meter}. The classifier
samples the target volume only at heights the map can know, and an occluder is classifiable
only where it rises above that plane at its range: trunks do, a chest-high obstacle from a
vehicle at head height does not. State it with Eq.~\eqref{eq:classify}.}

\subsection{Reappearance Belief and Planner Conditioning}
\label{sec:reappear}
Under \occluded{}, we propagate the predicted state across the occluder and
retain the set $\mathcal{R}$ of boundary crossings with arrival times, giving a
distribution over where and when the subject re-emerges. The viewpoint planner is
then conditioned on the loss class:

\begin{itemize}
  \item \occluded{}: hold a tight standoff, and bias the viewpoint toward the
        highest-probability element of $\mathcal{R}$ --- pre-position for the
        edge rather than chase the extrapolated point.
  \item \lost{}: abandon the standoff geometry and execute a search
        \tbd{define it: ascend-and-widen is the obvious choice and matches the
        loss response the safety layer already implements.}
\end{itemize}

\note{As built 2026-09-14 (docs/m6-loss-oracle.md, ``M6.4b''); every number below is a bench
figure from the real planner process on synthetic scans. Steps 3 and 4 are each behind their
own switch, off in code and on in the flight configuration, so the belief can be reported with
and without each --- that is the ablation. Nothing here has flown.}

\paragraph{Constructing $\mathcal{R}$} From the blocked ray nearest the target's
centre we sweep the bearing left and right in one-cell angular steps at the
target's range until \bench{three} consecutive columns are free, which gives the
occluder's two edges and its angular interval. We then roll $K=\bench{32}$
constant-turn-rate samples forward from the coasted state for \bench{\SI{3.2}{\second}}
at \bench{\SI{0.1}{\second}}, each with its own speed, turn rate and heading drawn about
the predictor's. A sample \emph{re-emerges} at its first point that lies outside the
occluder's interval and has a camera ray containing no \textsc{Occupied} cell
(\textsc{Unknown} is allowed: the space behind a trunk is unknown by nature).
Re-emergence points are clustered by side of the occluder; each cluster carries a
mean, a spread, an arrival-time distribution (min, median, max) and a mass, and
the mass that never re-emerges within the horizon is the \emph{still-hidden}
mass $w_h$. $\mathcal{R}$ is this set of clusters. The draws are seeded per event, so
a replay reproduces the belief.

\paragraph{Rollouts that go round} The samples walk on the same map. A step is
split into sub-steps no longer than a cell, and a position is blocked when the
walker's body --- sampled every half cell along its lateral diameter
($\pm\bench{\SI{0.25}{\meter}}$) and the radius ahead, at every height the map can
know --- touches an \textsc{Occupied} cell of the raw grid (inflation is sized for
the vehicle, not a person). A blocked sub-step tries headings
$\pm\bench{15^\circ}, \pm\bench{30^\circ}, \ldots, \pm\bench{90^\circ}$, the sample's own turn
side first, and takes the first free one \emph{for that sub-step only}; the
intended heading is tried again on the next, so the sample slides round the
obstacle and resumes its course. A coasted state that already lies inside an
obstacle is backed off along its reverse heading, and a cluster whose mean lands
in occupied cells reports its medoid. No new random draws are made, so with this
step disabled the belief is bit-for-bit the one above \bench{(36/36 events
identical, checked against the previous build)}. Without it, the exit test's
\bench{\SI{0.4}{\meter}} stop-short --- there so the subject's own still-occupied cells
never count as an occluder --- records a re-emergence \emph{on} the mapped face of a
trunk the sample's path crosses \bench{(\SI{0.03}{\meter} from an occupied cell)}; with
it, every re-emergence point keeps $\geq\bench{\SI{0.19}{\meter}}$ from every mapped cell
and the dominant side is unchanged.

\paragraph{What the map can vouch for} The map holds an obstacle as the shell of
the face the sensor has seen: a \SI{0.6}{\meter} trunk \SI{4}{\meter} away was
\bench{six} cells per height, its back half, interior and tangent zones
\textsc{Unknown}. The samples therefore keep out of the \emph{mapped} body of an
obstacle and out of nothing the map has never seen; a sample that gets round the
lit face's edge walks through the unseen back unopposed. A trunk mapped from two
sides is a closed shell, from one side a wall. This is a precondition of the
belief, general to any surface map, and belongs beside the raycasting
precondition of Sec.~\ref{sec:classify}.

\paragraph{Reduction to a goal and a yaw} The planner consumes a goal position
and a yaw target. The yaw target is the dominant cluster's mean. The goal is the
candidate on the standoff annulus of Sec.~\ref{sec:viewpoint} about that point
(\bench{24} azimuths $\times$ \bench{3} radii) that lies outside inflated occupancy, at least
$d_{\text{hard}}$ from the subject's coasted position, has a camera ray to the point
with no \textsc{Occupied} cell, and is nearest the vehicle. The tracker uses the
cluster's arrival-time window and its region as the re-acquisition gate
(Sec.~\ref{sec:goalcontract}).

\paragraph{Looking behind} When most samples stay hidden --- $w_h\geq\bench{0.5}$, or
no sample re-emerges at all --- pre-positioning for an edge is pointless. The yaw
target then becomes the hidden samples' mean position at $t=\bench{\SI{1.5}{\second}}$,
and the same annulus search runs about that point: its clear-ray requirement
forces the goal off the camera--occluder line, which is exactly the sidestep a
pilot makes to see behind a tree. $\mathcal{R}$ is left as it is. On the bench the
trigger is well separated: behind an r\,\bench{0.3}\,m trunk at \bench{\SI{1}{\meter\per\second}}
only \bench{0.31--0.59} of the samples stay hidden and the cluster goal already looks
behind the trunk; behind an r\,\bench{0.5}\,m trunk at \bench{\SI{0.7}{\meter\per\second}}
\bench{0.62--0.84} stay hidden and the look point fires on every event, the goal
\bench{\SI{1.0}{\meter}} off the camera--trunk line. \tbd{The trigger has no hysteresis; if
flight logs show the yaw target alternating when $w_h$ straddles the threshold,
add one and say so.}

\paragraph{Cost} With both steps enabled the full evaluation --- rays, extent,
rollouts, reduction --- takes \bench{\SI{2.5}{\milli\second}} at the 99th percentile on the
flight computer (\bench{\SI{0.4}{\milli\second}} without the obstacle-aware rollouts),
against an \SI{8}{\hertz} coast rate.

\subsection{Standoff Viewpoint Planning}
\label{sec:viewpoint}
The nominal (target-visible) behaviour places a goal on the annulus
$[d_{\min}, d_{\max}] = [\bench{3.0}, \bench{4.0}]$\,m about the subject with a
hard minimum standoff of \bench{\SI{2.5}{\meter}}, plus a separate yaw target so
the vehicle faces the subject rather than its velocity vector. Two choices matter
more than the geometry. First, the goal leads along the predicted path: a
point-to-point planner replans rest-to-rest as distance-to-goal approaches zero,
so a goal at the subject's current position induces continuous replanning exactly
when the vehicle is closest to the person. Second, yaw is a separate channel, and
if it is silently misrouted the vehicle flies a geometrically perfect standoff
while never pointing its camera at the subject, with no error reported anywhere.

\subsection{The Goal-Acknowledgment Contract}
\label{sec:goalcontract}
Safety-assured planners do not accept arbitrary goals. Ours~\cite{ren2025super}
may accept a requested goal, relocate it to the nearest free cell within a fixed
radius, or reject it and retain the previous one --- reporting none of these. For
a static waypoint this is benign; for a continuously moving standoff goal the
follower would believe the vehicle is heading somewhere it is not. We extend the
planner to preserve commanded altitude, emit an explicit
\textsc{Accepted}/\textsc{Snapped}/\textsc{Dropped} acknowledgment per request,
and \emph{echo the requested point} in it --- necessary because under
\textsc{Dropped} the reported active goal is the previous one, geometrically
unrelated to the request, so without the echo a subscriber cannot tell which
request an acknowledgment answers.

\subsection{Safety Layer}
\label{sec:safety}
Goal sources are arbitrated
$\textsc{Safety} \succ \textsc{Manual} \succ \textsc{Follow} \succ \textsc{Nav}$,
manual authority being enforced by the autopilot. Three properties define the
layer. \emph{One publishing path:} exactly one function may emit a goal, from a
single periodic timer, because a second path is a path that can bypass the veto.
\emph{The veto is independent of the follow product:} the human-proximity veto
(radius \bench{\SI{1.5}{\meter}}) subscribes to the raw detection stream --- every
person, every frame, including empty frames --- so a defect in viewpoint geometry
cannot defeat it. \emph{Silence is not a stop at the goal layer:} a stored goal
keeps the vehicle flying, so every stop condition \emph{publishes} a hold near
the current position rather than going quiet. Composed, these give the loss
response for free: each stage stops publishing when it loses its belief, and the
arbitration layer converts that silence into a commanded hold.

\note{Under the new thesis this is a secondary contribution. Keep it --- it is
what makes flying near the followed human defensible --- but it must not crowd
Sec.~\ref{sec:classify}. Cut to half a column if page-limited.}

%==============================================================================
\section{Implementation}
\label{sec:impl}

\paragraph{Process boundaries are a bandwidth decision} The organized point array
needed for per-detection ranging is $1280\times720\times3$ floats,
\bench{\SI{11.06}{\mega\byte}} per frame, \bench{\SI{88}{\mega\byte\per\second}}
at the \SI{8}{\hertz} detection rate. Our middleware could not carry it: \bench{a
\SI{2.76}{\mega\byte} message at \SI{30}{\hertz} already lost
${\approx}\SI{8}{\percent}$ of frames} in a Python subscriber, and best-effort
transport fragments an \SI{11}{\mega\byte} sample across thousands of datagrams
where one loss discards the sample. The array is built inside the perception
process and never serialized, with an off-by-default debug publisher for
visualization and replay sharing one deprojection implementation, so what is
recorded deprojects exactly as what is flown.
\note{Since 2026-09-13 the flown cloud is built on the GPU (cloud\_backend:
cuda) while the debug publisher stays numpy, so ``one implementation'' is now
``two implementations verified bit-identical'' (200/200 real frames; re-checked
by the node on its first frames every run). Reword, and add the CPU effect
(bag replay: $\sim$1.9 $\to$ 0.7 cores) once measured in flight.}

\paragraph{Failure detection} A watchdog monitors publication rates for every
sensor-derived stream, because a sensor can stop producing while its process, USB
enumeration and topic registrations all remain healthy \bench{(observed on the
depth camera)}.

\begin{table}[t]
\centering
\caption{Onboard compute. \bench{Blue} = bench measurement; in-flight pending.}
\label{tab:compute}
\footnotesize
\begin{tabular}{@{}lrr@{}}
\toprule
Process & CPU (\% of one core) & RSS (MB) \\
\midrule
LiDAR-inertial odometry   & \bench{39.6} & \bench{157} \\
Human perception          & \bench{60.1} & \bench{503} \\
Target manager            & \bench{4.0}  & \bench{59}  \\
Viewpoint planner         & \bench{0.7}  & \bench{25}  \\
Arbitration + veto        & \bench{10.4} & \bench{59}  \\
Loss classifier           & \tbd{}       & \tbd{}      \\
Trajectory planner        & \tbd{}       & \tbd{}      \\
Setpoint bridge           & \bench{34}   & \bench{69}  \\
\midrule
Total                     & \tbd{}       & \tbd{}      \\
\bottomrule
\end{tabular}
\end{table}

%==============================================================================
\section{Experiments}
\label{sec:results}

\note{Evaluation design follows the bibliography's Bucket 7. Do the offline
study FIRST --- it needs no flight time and de-risks the paper.}

\subsection{Offline: Prediction Under Missing Observations}
Before any flight, we evaluate the predictor in isolation on standard pedestrian
data with induced missing observations, following the protocol
of~\cite{chib2024trajimpute}. Metrics are ADE/FDE, reported both with complete
histories and under increasing missing-observation rates.

\begin{table}[t]
\centering
\caption{Prediction error under induced missing observations
(\cite{chib2024trajimpute} protocol). \tbd{all cells}}
\label{tab:prediction}
\footnotesize
\begin{tabular}{@{}lcccc@{}}
\toprule
& \multicolumn{2}{c}{Complete history} & \multicolumn{2}{c}{\tbd{X}\% missing} \\
\cmidrule(lr){2-3}\cmidrule(lr){4-5}
Predictor & ADE & FDE & ADE & FDE \\
\midrule
CTRV only                 & \tbd{} & \tbd{} & \tbd{} & \tbd{} \\
Learned only~\cite{bae2023eigentrajectory} & \tbd{} & \tbd{} & \tbd{} & \tbd{} \\
Fused (ours)              & \tbd{} & \tbd{} & \tbd{} & \tbd{} \\
\bottomrule
\end{tabular}
\end{table}

\subsection{Flight: Forest Trials}
\tbd{Venue, canopy density, trunk diameter distribution, subject walking speeds,
safety-pilot arrangement, geofence, abort criteria. Report trunk statistics ---
the occlusion argument is a claim about geometry and the reviewer needs the
geometry.}

\paragraph{Baselines} (1)~stock~\cite{ren2025super}: constant-velocity model with
a goal behind the target; (2)~our system with unconditional carry-forward and a
fixed missed-frame drop, i.e.\ ablating the classifier; (3)~a
re-acquisition-capable prior system~\cite{lin2026evatracker}; (4)~ours.
\note{Baseline 3 is the one reviewers will demand, and its code is public --- run
the real thing rather than a reimplementation if at all possible.}

\paragraph{Metrics}
\begin{itemize}
  \item Re-acquisition time after occlusion onset \note{the headline number}
  \item Tracking-loss duration, mean and max, per trial
  \item Visibility rate: fraction of mission time with the subject in view
  \item Tracking success rate over $N$ runs
  \item \textbf{Loss-classification accuracy}: \occluded{} vs.\ \lost{} against
        ground truth \tbd{ground-truth mechanism --- an RTK tag on the subject
        used for evaluation only, never for detection. Say this explicitly or a
        reviewer will assume you cheated.}
  \item Bearing error: angle between camera axis and subject
  \item Safety-pilot interventions per autonomous minute, with cause
\end{itemize}

\begin{table}[t]
\centering
\caption{Forest following performance over \tbd{N} trials per condition.}
\label{tab:following}
\footnotesize
\begin{tabular}{@{}lrrrr@{}}
\toprule
Method & Re-acq.\ (s) & Vis.\ rate (\%) & Success (\%) & Interv./min \\
\midrule
SUPER stock~\cite{ren2025super}    & \tbd{} & \tbd{} & \tbd{} & \tbd{} \\
Carry-forward (ablation)           & \tbd{} & \tbd{} & \tbd{} & \tbd{} \\
Eva-Tracker~\cite{lin2026evatracker} & \tbd{} & \tbd{} & \tbd{} & \tbd{} \\
\textbf{Ours}                      & \tbd{} & \tbd{} & \tbd{} & \tbd{} \\
\bottomrule
\end{tabular}
\end{table}

\subsection{Ablations}
\begin{enumerate}
  \item \textbf{With / without loss classification} --- \note{THE ablation. If
        this does not separate cleanly on re-acquisition time, the paper is an
        RA-L systems paper, not an ICRA algorithms paper. Run it first.}
  \item Learned vs.\ CTRV vs.\ fused prediction \emph{during} occlusion.
  \item Reappearance as a point vs.\ as a region.
  \item \tbd{ray count / threshold sensitivity for Eq.~\eqref{eq:classify}}
\end{enumerate}

\subsection{Safety-Layer Behaviour Under Induced Failure}
\tbd{Occlude the subject; kill perception mid-follow; introduce a second person
inside the re-acquisition band; starve the state estimate. Show the commanded
response is the designed one in each case. This subsection is cheap to run and
distinguishes the paper from tracking-accuracy-only submissions.}

%==============================================================================
\section{Lessons and Negative Results}
\label{sec:lessons}

\note{ICRA rewards this when specific and generalizable. Format each: symptom,
why it was invisible, what it implies. CUT THIS SECTION FIRST if at 6 pages.}

\paragraph{Latency no consumer detects} Our odometry produced correct values with
correct timestamps that were \bench{${\approx}\SI{1.4}{\second}$ old}. Every
downstream freshness check used message \emph{arrival} time, so nothing reported
a fault, and a stationary vehicle is indistinguishable either way. Cause: a
single-threaded executor whose always-overdue high-rate timer starved the sensor
callback, plus an input queue that filled at startup and never drained.
\bench{Corrected to \SIrange{49}{57}{\milli\second}.} \emph{Implication:} validate
freshness against the stamp, on a moving platform; and note that fixing latency
by shrinking a queue converts a queuing failure into a dropping failure that a
bench cannot exercise.

\paragraph{Occupancy maps cannot forget a person} Cells are freed only by
raycasting, disabled by default in the high-rate configuration, so a passing
person is inserted as permanent occupancy --- and a \emph{followed} person stamps
a wall behind themselves along the corridor the vehicle must fly.
\emph{Implication:} dynamic-obstacle handling is a precondition for aerial
following, not an enhancement, and it is a precondition for our own
classifier (Sec.~\ref{sec:classify}).

\paragraph{Frame conventions fail silently} Every frame error we made produced a
plausible, self-consistent, wrong system: a body origin that was the sensor's
internal IMU; an estimator delay double-counting a latency already in the
timestamp; a yaw convention differing from the obvious guess by $90^\circ$; a
mounting yaw unobservable from the estimator consuming it. None raised an error.
\emph{Implication:} assert against the \emph{plausible} wrong answer, not merely
against a wrong one.

\paragraph{Middleware determines architecture} Process boundaries were dictated
by transport bandwidth, not design preference. \emph{Implication:} measure the
transport before drawing the node diagram.

\paragraph{The companion link is flight-critical and unprofiled}
\bench{We recorded 22 autopilot hard faults across two firmware builds at a single
program counter in the companion-link client, traced to a use-after-free in the
serial driver's close/reopen path} --- i.e.\ triggered by the link loss most
likely to occur when something else has already gone wrong.
\note{2026-09-13: reproduced on demand on the bench; disabling RX DMA did not fix it; a
client-side change that keeps the UART open across session losses passed 30/30 induced
session losses but has not yet run the planned soak. If this paragraph stays, report the fix
only once the soak passes.}
\tbd{decide: paper, footnote, or cut. It is honest and unusual to report, and it
may read as off-topic for a tracking paper.}

%==============================================================================
\section{Limitations}
\label{sec:limitations}
\begin{itemize}
  \item \textbf{Single subject.} Group following and handoff are out of scope.
  \item \textbf{Re-identification is geometric.} Abstention under ambiguity is
        the correct conservative behaviour but ends follows an appearance model
        could continue. \tbd{quantify how often}
  \item \textbf{Classification depends on map quality} --- specifically on the
        map's ability to clear stale occupancy (Sec.~\ref{sec:classify}).
  \item \textbf{Active-stereo depth under canopy and sun} is uncharacterized.
  \item \textbf{Calibration is bench-solved}, not re-validated under
        motor-induced vibration.
  \item \textbf{No consent mechanism.} The system distinguishes the locked
        subject from bystanders and keeps distance from both, but ``the person
        the operator engaged on'' is not ``a person who agreed to be followed by
        an aircraft.'' We regard this as a real limitation of aerial following
        generally, and note that our proximity veto is the only mechanism in the
        system that protects a non-consenting bystander.
\end{itemize}

%==============================================================================
\section{Conclusion}
\label{sec:conclusion}
\tbd{Write last. Land: the geometry needed to know \emph{why} you lost someone is
already in the map you built to avoid hitting things, and using it turns one
undifferentiated failure into two with distinct, cheap, correct responses. Point
at released code.}

\section*{Acknowledgment}
\tbd{funding, lab, safety pilots, hardware loans}

%==============================================================================
\bibliographystyle{IEEEtran}
\bibliography{references}

\end{document}
